% ============================================================
%  AGRADECIMIENTOS
% ============================================================
\chapter*{Agradecimientos}
\addcontentsline{toc}{chapter}{Agradecimientos}
\enlargethispage{\baselineskip}

Al M. en C. Arturo Ronquillo Arvizu, director de esta tesis, por
haberme abierto las puertas del laboratorio, por la confianza de dejarme
intervenir un arreglo experimental en uso, y por la paciencia con que
acompañó un trabajo que avanzó a base de sesiones cortas y contratiempos
largos.

Al Laboratorio de Fotofísica y Películas Delgadas del Instituto de
Ciencias Aplicadas y Tecnología, por el espacio, el equipo y el tiempo de
instrumento sin los cuales este trabajo no habría sido posible.

Este trabajo recibió apoyo financiero del proyecto PAPIIT IN108024 y de
la Secretaría de Ciencia, Humanidades, Tecnología e Innovación (SECIHTI).

A la Dra. Ligia Pérez Cruz y a la Coordinación de Plataformas Oceanográficas
(COPO), por facilitar el equipo Minos X, empleado en el laboratorio para
la caracterización de muestras.

Al Ing. Francisco Ponce Núñez, a la Dra. María Guadalupe Ponce Vélez, a
la Q.F.B. Susana Santiago Pérez, al Ing. Santiago de Jesús Allende
Mendoza, al Lic. Carlos Daniel Gómez González, al Ing. Uriel Arana
Robles, a José Luis A. Ortiz Gordillo, al Dr. José Roberto Rivera
Hernández, al M. en C. Daniel Armando Pérez Calderón, al Dr. Armando
Reyes Montero, al Mtro. Fermín Sergio Castillo Sandoval, a Julio Antonio
Lestayo González, al Dr. Santiago Jesús Pérez Ruiz y al Mtro. Jorge
Arturo Real Pérez, por su valiosa colaboración y apoyo.

A Luis Pablo Rangel Mendieta, cuya celda fotoacústica e hidrófono son el
arreglo sobre el que se construyó esta automatización. Buena parte de lo
que aquí se reporta descansa en lo que él dejó hecho.

A la Universidad Nacional Autónoma de México y a la Facultad de
Ingeniería, por una formación que fue exigente en el mejor sentido de la
palabra.

A mis padres, a mis abuelos y a mis tíos y tías, por un apoyo que nunca
estuvo condicionado a que las cosas salieran bien. Lo que soy y lo que
pude terminar se lo debo a ellos.

A mis amigos, que me alentaron a seguir adelante y que estuvieron
presentes en los peores momentos. Terminar esto solo habría sido mucho
más difícil, y bastante menos llevadero.
